\documentclass[12pt]{article}
\usepackage{amsmath, amssymb}
\usepackage{geometry}
\geometry{margin=1in}
\usepackage{enumitem}

\title{Lecture Scribe\\
CSE 400: Fundamentals of Probability in Computing\\
Tutorial--2}
\date{}
\begin{document}

\maketitle

\textbf{Source: Tutorial-2 Document}

This scribe organizes and reconstructs the mathematical concepts, definitions, distributions, and general derivation techniques required to solve the problems in Tutorial--2. It is structured conceptually rather than question-wise and contains only material traceable to the provided document.

\section*{1. Continuous Uniform Distribution on an Interval}

\subsection*{1.1 Definition}

A continuous random variable $X$ is Uniform on $(a,b)$ if:

\[
f_X(x)=
\begin{cases}
\frac{1}{b-a}, & a<x<b \\
0, & \text{otherwise}
\end{cases}
\]

The cumulative distribution function (CDF):

\[
F_X(x)=
\begin{cases}
0, & x\le a \\
\frac{x-a}{b-a}, & a<x<b \\
1, & x\ge b
\end{cases}
\]

The expectation:

\[
E[X]=\frac{a+b}{2}
\]

\subsection*{1.2 Random Point on a Line Segment}

If a point is chosen at random on a line segment of length $L$, this means:

\[
X \sim \text{Uniform}(0,L)
\]

A geometric probability involving ratios (e.g., shorter/longer segment) requires:

\begin{itemize}
\item Expressing the ratio in terms of $X$
\item Identifying regions of the sample space where the inequality holds
\item Computing probability via integration over the appropriate interval
\end{itemize}

For continuous variables:

\[
P(A)=\int_{\text{region}} f_X(x)\,dx
\]

Since the density is constant for uniform distributions, probability reduces to:

\[
\frac{\text{length of favorable region}}{\text{total length}}
\]

\section*{2. Normal (Gaussian) Distribution and Bayesian Classification}

\subsection*{2.1 Definition}

A random variable $X\sim N(\mu,\sigma^2)$ has pdf:

\[
f_X(x)=\frac{1}{\sqrt{2\pi\sigma^2}}\exp\left(-\frac{(x-\mu)^2}{2\sigma^2}\right)
\]

Key parameters:

Mean: $\mu$

Variance: $\sigma^2$

\subsection*{2.2 Mixture Models and Prior Probabilities}

When a point is drawn from two possible populations:

With probability $\alpha$, it comes from one distribution.

With probability $1-\alpha$, from another.

This forms a mixture distribution:

\[
f(x)=\alpha f_1(x)+(1-\alpha)f_2(x)
\]

\subsection*{2.3 Posterior Probability (Bayes' Rule)}

If $A$ and $B$ partition the sample space:

\[
P(A\mid x)=\frac{f_A(x)P(A)}{f_A(x)P(A)+f_B(x)P(B)}
\]

Equal error probability condition implies balancing posterior probabilities.

Thus, one sets:

\[
P(A\mid x)=P(B\mid x)
\]

leading to an equation in prior probability $\alpha$.

\section*{3. Minimization of Expected Absolute Deviation}

We are interested in minimizing:

\[
E[|X-a|]
\]

\subsection*{3.1 General Form}

\[
E[|X-a|]=\int_{-\infty}^{\infty} |x-a|f_X(x)\,dx
\]

Split at $x=a$:

\[
= \int_{-\infty}^{a} (a-x)f_X(x)\,dx
+ \int_{a}^{\infty} (x-a)f_X(x)\,dx
\]

Differentiating with respect to $a$ gives:

\[
\frac{d}{da}E[|X-a|]=P(X\le a)-P(X\ge a)
\]

Setting derivative to zero:

\[
P(X\le a)=\frac{1}{2}
\]

Thus:

The minimizer of $E[|X-a|]$ is the median of $X$.

\subsection*{3.2 Uniform Case}

For $X\sim \text{Uniform}(0,A)$:

Median $= \frac{A}{2}$

\subsection*{3.3 Exponential Case}

If $X\sim \text{Exponential}(\lambda)$:

\[
f_X(x)=\lambda e^{-\lambda x}, \quad x>0
\]

The median $m$ satisfies:

\[
P(X\le m)=\frac{1}{2}
\]

\[
1-e^{-\lambda m}=\frac{1}{2}
\]

This determines the optimal location.

\section*{4. Exponential Distribution and Memoryless Property}

\subsection*{4.1 Definition}

If $X\sim \text{Exponential}(\lambda)$:

\[
f_X(x)=\lambda e^{-\lambda x}, \quad x>0
\]

\[
F_X(x)=1-e^{-\lambda x}
\]

\subsection*{4.2 Memoryless Property}

\[
P(X>t+s\mid X>t)=P(X>s)
\]

Proof:

\[
\frac{P(X>t+s)}{P(X>t)}=
\frac{e^{-\lambda(t+s)}}{e^{-\lambda t}}=e^{-\lambda s}
\]

\subsection*{4.3 Mixture of Exponentials}

If with probability $p_i$, rate is $\lambda_i$,

then survival probability:

\[
P(X>t)=\sum_i p_i e^{-\lambda_i t}
\]

Conditional survival:

\[
P(X>t+s\mid X>t)=
\frac{\sum_i p_i e^{-\lambda_i(t+s)}}
{\sum_i p_i e^{-\lambda_i t}}
\]

This is not memoryless unless rates are identical.

\section*{5. Poisson Distribution and Conditional Structure}

\subsection*{5.1 Definition}

$X\sim \text{Poisson}(\lambda)$

\[
P(X=k)=e^{-\lambda}\frac{\lambda^k}{k!}
\]

Mean and variance:

\[
E[X]=\lambda
\]

\subsection*{5.2 Sum of Independent Poissons}

If:

\[
X\sim \text{Poisson}(\lambda_1), \quad
Y\sim \text{Poisson}(\lambda_2)
\]

and independent, then:

\[
X+Y\sim \text{Poisson}(\lambda_1+\lambda_2)
\]

\subsection*{5.3 Conditional Distribution Given Total}

\[
P(X=k\mid X+Y=n)=
\frac{P(X=k,Y=n-k)}{P(X+Y=n)}
\]

Using independence:

\[
=\frac{
e^{-\lambda_1}\frac{\lambda_1^k}{k!}
e^{-\lambda_2}\frac{\lambda_2^{n-k}}{(n-k)!}
}
{
e^{-(\lambda_1+\lambda_2)}
\frac{(\lambda_1+\lambda_2)^n}{n!}
}
\]

This simplifies to a Binomial-type structure.

\section*{6. Transformation of Random Variables}

\subsection*{6.1 CDF Method}

If $Y=g(X)$:

\[
F_Y(y)=P(g(X)\le y)
\]

\[
f_Y(y)=\frac{d}{dy}F_Y(y)
\]

\subsection*{6.2 Piecewise Transformations}

If $g$ is not one-to-one:

\[
f_Y(y)=\sum_{x:g(x)=y}
\frac{f_X(x)}{|g'(x)|}
\]

Valid when inverse branches exist.

\section*{7. Binomial Approximation via Independent Movements}

If each period:

Probability $p$ of upward movement

Independent across periods

Then after $n$ periods:

\[
X\sim \text{Binomial}(n,p)
\]

\[
P(X=k)=\binom{n}{k}p^k(1-p)^{n-k}
\]

Large $n$ approximation:

\[
X\approx N(np,np(1-p))
\]

Probability involving percentage increase translates into inequality in $X$.

\section*{8. Conditional Probability with Exponential Variables}

Using survival function:

\[
P(X>t)=e^{-\lambda t}
\]

Conditional probability:

\[
P(X\ge t+s\mid X\ge t)
=\frac{P(X\ge t+s)}{P(X\ge t)}
\]

Apply exponential survival form.

\section*{9. Joint Transformation to Radial Variable}

Let $X$ and $Y$ independent with pdfs $f_X(x), f_Y(y)$.

Define:

\[
Z=X^2+Y^2
\]

To obtain pdf:

Use joint density:

\[
f_{X,Y}(x,y)=f_X(x)f_Y(y)
\]

Transform to polar coordinates:

\[
x=r\cos\theta, \quad y=r\sin\theta
\]

Jacobian:

\[
J=r
\]

Integrate over angle:

\[
f_Z(r)=\int_0^{2\pi}
f_X(r\cos\theta)f_Y(r\sin\theta) r\, d\theta
\]

\section*{10. Order Statistics and Joint Distributions}

Let $X_1,X_2,\dots$ be i.i.d. with CDF $F$.

Define:

\[
M_n=\max(X_1,\dots,X_n)
\]

\subsection*{10.1 Distribution of Maximum}

\[
P(M_n\le x)=P(X_1\le x,\dots,X_n\le x)
\]

Using independence:

\[
= [F(x)]^n
\]

Thus:

\[
F_{M_n}(x)=[F(x)]^n
\]

\subsection*{10.2 Joint Distribution of $M_n$ and $M_{n+1}$}

Observe:

\[
M_{n+1}=\max(M_n,X_{n+1})
\]

Joint event:

\[
P(M_n\le x, M_{n+1}\le y)
\]

Case distinction:

If $y<x$: probability is zero.

If $y\ge x$:

\[
P(M_n\le x, X_{n+1}\le y)
=[F(x)]^n F(y)
\]

This characterizes the joint CDF structure.

\section*{Summary of Core Concepts}

This tutorial requires mastery of:

\begin{itemize}
\item Uniform distribution and geometric probability
\item Normal distributions and mixture models
\item Bayes' rule for posterior probabilities
\item Minimization of expected absolute deviation (median property)
\item Exponential distribution and memorylessness
\item Mixture of exponentials
\item Poisson distribution and conditional structure
\item Binomial modeling of repeated independent trials
\item Normal approximation to Binomial
\item Transformation of random variables (CDF and Jacobian methods)
\item Order statistics and maxima distributions
\end{itemize}

\end{document}
